\chapter{Разработка алгоритма с квантованием}
\label{ch:chap3}

В третьем разделе работы базовый алгоритм из Главы~\ref{ch:chap2} расширен модулем векторного квантования визуальных представлений. Цель такого расширения двойная: с одной стороны, дискретизация формирует информационное «бутылочное горлышко», подавляющее нерелевантные доменные компоненты входа \cite{yang2023vector} и тем самым повышающее устойчивость политики; с другой стороны, совместное обучение декодера изображения по полученным дискретным кодам обеспечивает дополнительный обучающий сигнал и способствует формированию более информативных признаков. Согласно техническому заданию, итоговый алгоритм должен обеспечивать успешность не менее 75\,\% эпизодов на наборе robomimic~v0.1; ниже описана архитектура, обоснование принятых проектных решений и схема обучения.

\section{Концептуальная схема}
\label{sec:vq-overview}

Для каждой камеры реализуется единая последовательная схема обработки кадра: изображение проходит DINOv2-энкодер с обучаемой проекционной головой; получаемая пространственная карта признаков пропускается через пространственный векторный квантователь с EMA-обновлением кодовой книги; квантованная карта параллельно поступает на декодер изображения, обеспечивающий пиксельную реконструкцию, и -- после линейной проекции -- на трансформерную политику, идентичную описанной в Главе~\ref{ch:chap2}. Полная функция потерь объединяет три слагаемых:
\begin{equation}
    \label{eq:vq-total-loss}
    \mathcal{L} = \mathcal{L}_{\text{policy}} + \beta \,\mathcal{L}_{\text{VQ}}
                  + \lambda \,\mathcal{L}_{\text{IR}},
\end{equation}
где $\mathcal{L}_{\text{policy}}$ -- та же отрицательная логарифмическая правдоподобность GMM-головы, что и в~\eqref{eq:bc-loss}, $\mathcal{L}_{\text{VQ}}$ -- штраф за расхождение энкодера с выбранным кодовым вектором (commitment loss), а $\mathcal{L}_{\text{IR}}$ -- ошибка реконструкции изображения. Весовые коэффициенты $\beta$ и $\lambda$ задаются как гиперпараметры (в реализации они равны
единице). 

\section{Адаптация поблочного квантования к DINOv2}
\label{sec:vq-block}

Идея пространственно-блочного векторного квантования восходит к VQ-VAE \cite{vandeoord2017neural}: сверточный энкодер выдает карту признаков $H{\times}W{\times}D$, и каждая из $H{\cdot}W$ позиций квантуется независимо ближайшим вектором кодовой книги. В работе Порихиса и соавторов \cite{porichis2024imitation} такая схема применена к небольшой сверточной сети для задачи робастного захвата плодов: пространственная локальность выходов свертки обеспечивает прямую связь каждого квантуемого блока с областью исходного изображения (receptive field), что и делает квантование осмысленным.

Замена базовой зрительной сети на DINOv2 не противоречит этой логике, а, наоборот, делает связь с исходным изображением еще более явной. Патч-эмбеддинг DINOv2 разбивает изображение $224{\times}224$ на регулярную сетку патчей $14{\times}14$, поэтому каждому патч-токену соответствует четко определенная область $14{\times}14$~пикселей входного кадра. После прогона через все блоки трансформера и финальную нормализацию патч-токены сохраняют пространственную организацию: их можно преобразовать в карту $[B, 384, 16, 16]$, где позиция $(i,j)$ напрямую соответствует одной и той же области изображения, что и при сверточной обработке. Тем самым поблочная схема \cite{vandeoord2017neural, esser2021taming, porichis2024imitation} переносится на ViT-энкодер без концептуальных изменений, и каждая из 256 ячеек $16{\times}16$ карты квантуется независимо. Альтернативные варианты -- глобальный пулинг или проекция всей карты в один вектор -- разрушали бы пространственную структуру и лишали бы дискретные коды интерпретации в терминах исходного кадра, поэтому от них отказались.

\section{Проекционный слой и сжатие размерности}
\label{sec:vq-proj}

Прямое квантование 384-мерных патч-токенов DINOv2 нецелесообразно по двум причинам. Во-первых, кодовая книга размерности 384 существенно увеличивает число параметров и время поиска ближайшего вектора. Во-вторых, DINOv2 обучается без учителя на изображениях из ImageNet-подобной коллекции, и большинство компонент его выхода нерелевантно симулированной сцене манипулятора. Для согласования размерностей и подавления избыточных степеней свободы добавлен обучаемый проекционный слой: свертка $1{\times}1$ с входной размерностью 384 и выходной $D = 64$. Поскольку свертка $1{\times}1$ применяется поточечно ко всем позициям карты признаков, она не нарушает пространственного соответствия патч-токенов и областей кадра. Полученная карта $[B, 64, 16, 16]$ и составляет вход в квантующий модуль.

Помимо проекционной свертки, обучаемыми остаются последний блок DINOv2 и финальная нормализация, как и в Главе~\ref{ch:chap2}. Все остальные параметры энкодера заморожены, что сохраняет универсальные свойства предобученных представлений.

\section{Пространственный векторный квантователь с EMA}
\label{sec:vq-quantizer}

Пусть $e_t \in \mathbb{R}^{D \times H \times W}$ -- выход проекционного слоя для одного кадра. Кодовая книга задается таблицей $E = \{b_i\}_{i=1}^{N}$, $b_i \in \mathbb{R}^{D}$, $N = 1024$. Каждая позиция $(h, w)$ карты признаков $e_t$ независимо квантуется ближайшим вектором кодовой книги:
\begin{equation}
    \label{eq:vq-argmin}
    k_{t}(h, w) = \arg\min_{i \in \{1, \dots, N\}} \,\bigl\| e_{t}(h, w) - b_{i} \bigr\|_{2}^{2},
\end{equation}
после чего составляется квантованный тензор $z_t(h, w) = b_{k_t(h,w)}$. Для эффективности расчет расстояний выполняется по разложению $\|e - b\|^{2} = \|e\|^{2} + \|b\|^{2} - 2\,\langle e, b\rangle$.

Поскольку операция $\arg\min$ недифференцируема, для прохождения градиентов от политики и реконструкции в энкодер используется straight-through-estimator
\cite{vandeoord2017neural}:
\begin{equation}
    \label{eq:vq-ste}
    z_t \;=\; e_t \;+\; \mathrm{sg}\!\left(z_t - e_t\right),
\end{equation}
где $\mathrm{sg}(\cdot)$ -- оператор остановки градиента. На прямом проходе результат равен квантованному значению, на обратном -- эквивалентен тождественному преобразованию.

Кодовая книга $E$ обновляется методом экспоненциального скользящего среднего, а не градиентным спуском, что улучшает устойчивость обучения и согласуется с реализацией \cite{vandeoord2017neural}. Обозначим через $r_i^{(t)}$ -- среднее по мини-батчу значение тех векторов $e_t(h,w)$, которые были квантованы в код~$i$, и через $n_i^{(t)}$ -- число таких векторов. Поддерживаются EMA-статистики
\begin{equation}
    \label{eq:vq-ema}
    \begin{aligned}
        N_i &\leftarrow \gamma\, N_i + (1 - \gamma)\, n_i^{(t)},\\
        m_i &\leftarrow \gamma\, m_i + (1 - \gamma)\, n_i^{(t)} r_i^{(t)},\\
        b_i &\leftarrow m_i \,/\, \tilde{N}_i,
    \end{aligned}
\end{equation}
где $\gamma = 0{,}99$ -- коэффициент EMA, а $\tilde{N}_i$ -- сглаженное по Лапласу значение $N_i$, исключающее вырождение редко используемых кодов. Для контроля разнообразия используется метрика перплексии $\exp\!\bigl(-\sum_i p_i \log p_i\bigr)$, где $p_i$ -- доля позиций, попавших в
код~$i$ за мини-батч; перплексия фиксируется в процессе обучения.

Дополнительно к этому энкодер регуляризуется commitment-потерей, гарантирующей, что выход энкодера не уходит произвольно далеко от ближайшего кодового вектора:
\begin{equation}
    \label{eq:vq-commit}
    \mathcal{L}_{\text{VQ}} = \beta_0 \,\bigl\| \mathrm{sg}(z_t) - e_t \bigr\|_{2}^{2},
\end{equation}
где $\beta_0 = 0{,}25$. Эта потеря, в отличие от EMA, действует именно на градиент энкодера и не затрагивает кодовую книгу.

\section{Декодер изображения и потеря реконструкции}
\label{sec:vq-decoder}

Параллельно политике обучается небольшой сверточный декодер изображения, восстанавливающий исходный кадр по квантованной карте $z_t$. Декодер состоит из четырех блоков \texttt{ConvTranspose2d} с шагом~2 (последовательность пространственных размеров $16 \to 32 \to 64 \to 128 \to 256$), каждый с BatchNorm и ReLU, и финальной сверткой $3{\times}3$ с сигмоидной активацией. Выход при необходимости интерполируется до целевого размера $224{\times}224$. Используется попиксельная бинарная кросс-энтропия:
\begin{equation}
    \label{eq:vq-recon}
    \mathcal{L}_{\text{IR}} = -\frac{1}{C \cdot H_o \cdot W_o}
       \sum_{c, h, w}
       \Bigl[
         I^{\ast}_{c,h,w} \log \hat{I}_{c,h,w}
         + (1 - I^{\ast}_{c,h,w}) \log (1 - \hat{I}_{c,h,w})
       \Bigr],
\end{equation}
где $\hat{I}$ -- выход декодера, $I^{\ast}$ -- целевое изображение, приведенное к размеру выхода декодера билинейной интерполяцией. Перед вычислением логарифмов значения ограничиваются диапазоном $[10^{-6}, 1 - 10^{-6}]$, что обеспечивает численную устойчивость при использовании случайных кадрировок.

Введение реконструкционной ветви дает дополнительный обучающий сигнал: чтобы по дискретным кодам $z_t$ оказалось возможным восстановить изображение, кодовая книга вынуждена сохранять достаточно информации о форме и расположении объектов, а не только о признаках, полезных непосредственно для классификации действия. Это особенно важно в режиме малого числа демонстраций, когда сигнал от $\mathcal{L}_{\text{policy}}$ может быть слишком слабым для качественного обучения зрительной части.

\section{Сопряжение с трансформерной политикой}
\label{sec:vq-policy}

После квантования карта $z_t \in \mathbb{R}^{64 \times 16 \times 16}$ разворачивается в одномерный вектор и пропускается через линейный слой размерности 128, формируя представление одной камеры на шаге $t$. Представления двух камер конкатенируются с низкоразмерным состоянием робота, образуя итоговый эмбеддинг наблюдения $h_t$. Дальнейшая обработка полностью совпадает с Главой~\ref{ch:chap2}: последовательность $h_{t-T+1:t}$ длиной $T = 10$ подается в декодерный трансформер с 6 слоями, эмбеддингом 512 и 8 головами внимания, поверх которого работает GMM-голова с пятью компонентами.

Принципиально важным проектным решением является то, что между квантующим модулем и трансформером сохраняется обычная дифференцируемая связь через straight-through-estimator. Это означает, что градиент политики достигает энкодера и проекционного слоя так же, как и в базовой версии, и кодовая книга не превращается в обучаемое узкое место.

\section{Поток градиентов и совместная оптимизация}
\label{sec:vq-grad}

Распределение градиентов по компонентам системы устроено согласно \cite{vandeoord2017neural,porichis2024imitation}:
\begin{itemize}
    \item потеря $\mathcal{L}_{\text{policy}}$ через straight-through-estimator \eqref{eq:vq-ste} проникает в проекционный слой и обучаемые компоненты DINOv2;
    \item потеря $\mathcal{L}_{\text{IR}}$ распространяется по той же траектории, дополнительно обновляя сверточный декодер;
    \item commitment-потеря $\mathcal{L}_{\text{VQ}}$ влияет только на энкодер: благодаря $\mathrm{sg}$ в~\eqref{eq:vq-commit} она не вносит вклада в обновление кодовой книги;
    \item кодовая книга $E$ обновляется только статистиками EMA из~\eqref{eq:vq-ema} и не имеет собственных градиентов.
\end{itemize}
Такое разделение ролей обеспечивает стабильность совместной оптимизации трех потерь. Соответствующая программная реализация выполнена в виде модуля \texttt{VisualCoreVQVAE}, наследующего интерфейс \texttt{VisualCore} библиотеки robomimic, и алгоритма \texttt{BC\_Transformer\_GMM\_VQVAE}, расширяющего штатный \texttt{BC\_Transformer\_GMM} сбором $\mathcal{L}_{\text{VQ}}$ и $\mathcal{L}_{\text{IR}}$ со всех зрительных энкодеров и взвешенным их добавлением к $\mathcal{L}_{\text{policy}}$ согласно~\eqref{eq:vq-total-loss}.

Гиперпараметры обучения совпадают с базовой конфигурацией Главы~\ref{ch:chap2} за двумя исключениями: размер мини-батча уменьшен до 8 (что обусловлено дополнительными требованиями декодера и квантователя к памяти GPU), а в качестве зрительного энкодера выбран \texttt{VisualCoreVQVAE} с описанными выше параметрами квантователя ($N = 1024$, $D = 64$, $\gamma = 0{,}99$, $\beta_0 = 0{,}25$). Прочие настройки -- скорость обучения, число эпох, длина окна, размер контекста трансформера, состав наблюдений и параметры случайной кадрировки -- не менялись.

\section{Результаты}
\label{sec:vq-results}

Обучение с описанной конфигурацией выполнено на подмножестве Proficient-Human набора robomimic~v0.1. Оценка политики проводилась в исходной среде robosuite путем развертывания: каждые 10 эпох запускалось 10 тестовых эпизодов с горизонтом 400 шагов. При тестировании алгоритма с пространственно-блочным квантованием визуальных представлений на задаче Lift из robomimic~v0.1 был достигнут уровень успешности $\text{Success Rate} = 0{,}96$, что превышает требуемый порог в 0{,}75 и подтверждает работоспособность реализации. Дальнейшее систематическое исследование влияния модуля квантования на робастность политики к доменным сдвигам выходит за рамки настоящей работы.

\endinput
